% ============================================================
% PRESENTACIÓN DE LA MATERIA
% ============================================================
\section*{PRESENTACIÓN DE LA MATERIA}
\subsection*{Seguridad Operacional y Gestión en Operaciones Aéreas}

\begin{center}
\small
\begin{tabular}{|p{0.32\linewidth}|p{0.58\linewidth}|}
\hline
\textbf{Escuela} & E.E.S.T. N° 7 ``Taller Regional Quilmes'' (IMPA) \\ \hline
\textbf{Especialidad} & Técnico/a en Aviónica \\ \hline
\textbf{Año} & 7mo año \\ \hline
\textbf{Campo} & Formación Científico Tecnológica \\ \hline
\textbf{Módulo} & Seguridad Operacional y Gestión en Operaciones Aéreas \\ \hline
\textbf{Entorno} & Legislación Aeronáutica \\ \hline
\textbf{Carga horaria} & 4 horas semanales -- 144 horas anuales \\ \hline
\textbf{Marco curricular} & Diseño Curricular 2025 de la Educación Secundaria Técnica, Provincia de
Buenos Aires \\ \hline
\textbf{Docente} & Téc. Aviónico Julián Scicolone -- scicolonejulian@impatrq.com \\ \hline
\end{tabular}
\end{center}

% ------------------------------------------------------------
\section*{FUNDAMENTACIÓN}

Un técnico aviónico no solo repara equipos: trabaja dentro de un sistema regulado, firma
documentos con valor legal y forma parte de organizaciones donde un error puede costar vidas. Por
eso el Diseño Curricular propone en este módulo el análisis integral de la \textbf{gestión de los
recursos humanos en organizaciones aéreas}, reconociendo la actuación humana como un factor
crítico en la operación y el mantenimiento, el estudio del \textbf{marco legal vigente} (Código
Aeronáutico, anexos de la OACI y regulaciones del Estado que rigen las licencias), el
\textbf{rol del técnico como representante técnico} de los talleres y la \textbf{oficina técnica}
como ámbito profesional, y la \textbf{prevención de accidentes} mediante las mejores prácticas
de gestión de la seguridad operacional.

\begin{recuadro}
\textbf{Un dato que le da sentido a esta materia.} Desde la Resolución ANAC N° 443/2025, el
\textbf{título de Técnico en Aviónica emitido por una escuela técnica} es uno de los títulos
aceptados para solicitar la \textbf{licencia de Mecánico de Equipos Radioeléctricos de Aeronaves
(MERA)}. Todo lo que se estudia en este manual (las regulaciones, la documentación, los factores
humanos y la seguridad operacional) es lo que se espera de un titular de esa licencia.
\end{recuadro}

% ------------------------------------------------------------
\section*{OBJETIVOS DE APRENDIZAJE}
Según el Diseño Curricular 2025, al finalizar el módulo se espera que las y los estudiantes
puedan:
\begin{itemize}
  \item Utilizar conceptos, herramientas, métodos y técnicas de la \textbf{gestión
        organizacional} de una empresa del ámbito aeronáutico en cualquiera de sus formas.
  \item Utilizar conceptos, herramientas, métodos y técnicas de la \textbf{gestión de calidad e
        impacto ambiental}.
  \item Utilizar conceptos, herramientas, métodos y técnicas de la \textbf{gestión de los
        recursos humanos} a través de programas de CRM o MRM.
  \item Interrelacionar la gestión global y la gestión de sus actividades específicas.
  \item Comprender el proceso productivo y la estructura organizacional de una empresa del sector
        aeronáutico (comercialización, selección y abastecimiento de sistemas, componentes y
        partes), en el marco de la normativa aeronáutica nacional e internacional vigente.
  \item Planificar, programar y organizar la producción, la prestación de servicios y el
        mantenimiento.
  \item \textbf{Planificar y gestionar un SMS} en un taller u organización aeronáutica.
  \item Identificar proyectos productivos y de servicio, evaluar su factibilidad
        técnico-económica, programar su implementación y su gestión.
\end{itemize}

% ------------------------------------------------------------
\section*{ORGANIZACIÓN DEL MANUAL}

\begin{center}
\small
\begin{tabular}{|c|>{\raggedright\arraybackslash}p{0.48\linewidth}|>{\raggedright\arraybackslash}p{0.33\linewidth}|}
\hline
\textbf{Unidad} & \textbf{Contenidos} & \textbf{Bloque del Diseño Curricular} \\ \hline
1 & Derecho aeronáutico: OACI, Código Aeronáutico, ANAC, JST, RAAC. Rol regulador del Estado. &
Derecho aeronáutico y aeronavegabilidad \\ \hline
2 & Aeronavegabilidad y certificación: aeronaves, partes, documentación, ATA 100/104, normas ISO,
AN, MS, NAS, MIL. & Derecho aeronáutico y aeronavegabilidad \\ \hline
3 & Licencias del personal (RAAC 65), representante técnico y oficina técnica. & Derecho
aeronáutico y aeronavegabilidad; Oficina técnica \\ \hline
4 & Organizaciones aeronáuticas: talleres, operadores, planificación del mantenimiento,
abastecimiento, control de gestión, economía del operador, calidad, ambiente y salud laboral. &
Organizaciones aeronáuticas \\ \hline
5 & Actuación humana: factores humanos, error humano, MRM, MEDA, reportes, primeros auxilios. &
Actuación humana \\ \hline
6 & Prevención de accidentes: riesgo, causalidad, investigación, SMS, costos de la seguridad
operacional, casos reales. & Prevención de accidentes \\ \hline
7 & Operaciones aéreas: tipos de operación, espacio aéreo y cartas aeronáuticas. &
Complementario \\ \hline
PI & Proyecto integrador: SMS del taller de aviónica de la escuela. & Prácticas formativas \\
\hline
\end{tabular}
\end{center}

\paragraph*{Metodología.} Cada unidad combina explicación teórica, lectura de normativa real,
análisis de casos (accidentes e incidentes investigados) y prácticas formativas en el taller.
Al final de cada unidad hay actividades. El año cierra con un \textbf{proyecto integrador}: el
diseño e implementación de un Sistema de Gestión de la Seguridad Operacional (SMS) para el
taller de la escuela.

\paragraph*{Fuentes.} Las regulaciones aeronáuticas cambian con frecuencia. Este manual está
actualizado a noviembre de 2026; antes de usar un dato normativo en un trabajo real, verificar
siempre la versión vigente en el sitio de la ANAC (\texttt{argentina.gob.ar/anac}) y en el
Boletín Oficial.

% ------------------------------------------------------------
\clearpage
\section*{HOJA DE RUTA DEL ESTUDIANTE}
\subsection*{Lo que vas a tener hecho al terminar el año}

\vspace{2pt}
{\small
\renewcommand{\arraystretch}{1.3}
\begin{tabular}{|c|>{\raggedright\arraybackslash}p{0.54\linewidth}|>{\centering\arraybackslash}p{0.08\linewidth}|>{\centering\arraybackslash}p{0.17\linewidth}|}
\hline
\textbf{OK} & \textbf{Logro} & \textbf{Unidad} & \textbf{Firma docente} \\ \hline
\casilla & Esquema institucional del sistema aeronáutico argentino (OACI, ANAC, JST, EANA, ORSNA) &
1 & \\ \hline
\casilla & Cuadro comparativo de normativa nacional e internacional & 1 & \\ \hline
\casilla & Verificación simulada de la aeronavegabilidad de una aeronave (documentos a bordo,
AD, registros) & 2 & \\ \hline
\casilla & Búsqueda de una tarea en un manual por capítulo ATA y de una parte en el IPC & 2 & \\
\hline
\casilla & Tabla comparativa de normas ISO, AN, MS, NAS y MIL & 2 & \\ \hline
\casilla & Carpeta de solicitud simulada de la licencia MERA & 3 & \\ \hline
\casilla & Organigrama funcional de un TAR con el rol del representante técnico & 3 & \\ \hline
\casilla & Registro de mantenimiento y retorno al servicio completados & 3 & \\ \hline
\casilla & Plan de mantenimiento e indicadores de gestión de una flota simulada & 4 & \\ \hline
\casilla & Caso práctico de economía del operador (costo de mantenimiento por hora de vuelo) & 4
& \\ \hline
\casilla & Simulación de pase de turno con planilla MRM & 5 & \\ \hline
\casilla & Reporte e investigación de un error con el método MEDA & 5 & \\ \hline
\casilla & Práctica de primeros auxilios (RCP y control de hemorragias) & 5 & \\ \hline
\casilla & Análisis de un accidente real con el modelo de Reason & 6 & \\ \hline
\casilla & Matriz de riesgos del taller de la escuela & 6 & \\ \hline
\casilla & Informe que justifique una inversión en prevención & 6 & \\ \hline
\casilla & Lectura de una carta aeronáutica y de un NOTAM & 7 & \\ \hline
\casilla & Proyecto integrador: manual SMS del taller presentado & PI & \\ \hline
\end{tabular}
}

\vspace{8pt}
\noindent Estudiante: \dotfill \quad Curso: \makebox[3cm]{\dotfill}

\vspace{6pt}
\noindent Logros completados: \makebox[1.5cm]{\dotfill} de 18 \hfill Fecha: \makebox[4cm]{\dotfill}
